% GENERATED FILE. Edit canonical lessons, then run npm run generate:books.
% canonical-source-hash: fnv1a64:35c7cfe89f8508f7
% canonical-lessons: TE-C220-ucitam, TE-C220-nindu, TE-C220-lavu, TE-C220-sannani, TE-C220-pacci

\chapter{Free, Full, Fat, Thin, Raw}
\label{ch:te-free-full-fat-thin-raw}
\hlchaptermodality{220}

\begin{chapteropening}
\textbf{By the end of this chapter:} \emph{I can say free, full, fat, thin and raw.}

Complete the last lesson of chapter 220: I can say free, full, fat, thin and raw.
\end{chapteropening}

\section[\texorpdfstring{\te{ఉచితం} (ucitaṁ)}{ఉచితం (ucitaṁ)}]{\te{ఉచితం} (ucitaṁ) — free (of charge)}

\label{lesson:TE-C220-ucitam}



\begin{quote}
Before the new one: say the Telugu for a development, then the Telugu for motivation.
\end{quote}

\subsection*{You'll want to know: \te{ఉచితం}}
\textbf{\te{ఉచితం}} — \emph{ucitaṁ} — \textquotedblleft{}free (of charge)\textquotedblright{}.

\begin{grammarlens}[title={What you've built}]
One more word for this chapter.
\end{grammarlens}

\subsection*{Your turn}
\begin{itemize}
\raggedright
  \item \emph{Say it:} \emph{ucitaṁ}
  \item \emph{Say it:} \emph{ucitaṁ}, once more
  \item \emph{Say it:} say \emph{ucitaṁ}
  \item \emph{Recall:} say the Telugu for a development, then the Telugu for motivation, then say \emph{ucitaṁ} again
\end{itemize}

\begin{tcolorbox}[breakable,colback=teal!4,colframe=teal!35!black,title={Before you move on}]
What does \te{ఉచితం} mean? (\textquotedblleft{}Free (of charge)\textquotedblright{}.) Say it once more.
\end{tcolorbox}
\section[\texorpdfstring{\te{నిండు} (niṇḍu)}{నిండు (niṇḍu)}]{\te{నిండు} (niṇḍu) — full}

\label{lesson:TE-C220-nindu}



\begin{quote}
Before the new one: say the Telugu for motivation, then the Telugu for free.
\end{quote}

\subsection*{You'll want to know: \te{నిండు}}
\textbf{\te{నిండు}} — \emph{niṇḍu} — \textquotedblleft{}full\textquotedblright{}.

\begin{grammarlens}[title={What you've built}]
One more word for this chapter.
\end{grammarlens}

\subsection*{Your turn}
\begin{itemize}
\raggedright
  \item \emph{Say it:} \emph{niṇḍu}
  \item \emph{Say it:} \emph{niṇḍu}, once more
  \item \emph{Say it:} say \emph{niṇḍu}
  \item \emph{Recall:} say the Telugu for motivation, then the Telugu for free, then say \emph{niṇḍu} again
\end{itemize}

\begin{tcolorbox}[breakable,colback=teal!4,colframe=teal!35!black,title={Before you move on}]
What does \te{నిండు} mean? (\textquotedblleft{}Full\textquotedblright{}.) Say it once more.
\end{tcolorbox}
\section[\texorpdfstring{\te{లావు} (lāvu)}{లావు (lāvu)}]{\te{లావు} (lāvu) — fat, thick}

\label{lesson:TE-C220-lavu}



\begin{quote}
Before the new one: say the Telugu for free, then the Telugu for full.
\end{quote}

\subsection*{You'll want to know: \te{లావు}}
\textbf{\te{లావు}} — \emph{lāvu} — \textquotedblleft{}fat, thick\textquotedblright{}.

\begin{grammarlens}[title={What you've built}]
One more word for this chapter.
\end{grammarlens}

\subsection*{Your turn}
\begin{itemize}
\raggedright
  \item \emph{Say it:} \emph{lāvu}
  \item \emph{Say it:} \emph{lāvu}, once more
  \item \emph{Say it:} say \emph{lāvu}
  \item \emph{Recall:} say the Telugu for free, then the Telugu for full, then say \emph{lāvu} again
\end{itemize}

\begin{tcolorbox}[breakable,colback=teal!4,colframe=teal!35!black,title={Before you move on}]
What does \te{లావు} mean? (\textquotedblleft{}Fat, thick\textquotedblright{}.) Say it once more.
\end{tcolorbox}
\section[\texorpdfstring{\te{సన్నని} (sannani)}{సన్నని (sannani)}]{\te{సన్నని} (sannani) — thin, slim}

\label{lesson:TE-C220-sannani}



\begin{quote}
Before the new one: say the Telugu for full, then the Telugu for fat.
\end{quote}

\subsection*{You'll want to know: \te{సన్నని}}
\textbf{\te{సన్నని}} — \emph{sannani} — \textquotedblleft{}thin, slim\textquotedblright{}.

\begin{grammarlens}[title={What you've built}]
One more word for this chapter.
\end{grammarlens}

\subsection*{Your turn}
\begin{itemize}
\raggedright
  \item \emph{Say it:} \emph{sannani}
  \item \emph{Say it:} \emph{sannani}, once more
  \item \emph{Say it:} say \emph{sannani}
  \item \emph{Recall:} say the Telugu for full, then the Telugu for fat, then say \emph{sannani} again
\end{itemize}

\begin{tcolorbox}[breakable,colback=teal!4,colframe=teal!35!black,title={Before you move on}]
What does \te{సన్నని} mean? (\textquotedblleft{}Thin, slim\textquotedblright{}.) Say it once more.
\end{tcolorbox}
\section[\texorpdfstring{\te{పచ్చి} (pacci)}{పచ్చి (pacci)}]{\te{పచ్చి} (pacci) — raw, unripe}

\label{lesson:TE-C220-pacci}



\begin{quote}
Before the new one: say the Telugu for fat, then the Telugu for thin.
\end{quote}

\subsection*{You'll want to know: \te{పచ్చి}}
\textbf{\te{పచ్చి}} — \emph{pacci} — \textquotedblleft{}raw, unripe\textquotedblright{}.

\begin{grammarlens}[title={What you've built}]
One more word for this chapter.
\end{grammarlens}

\subsection*{Your turn}
\begin{itemize}
\raggedright
  \item \emph{Say it:} \emph{pacci}
  \item \emph{Say it:} \emph{pacci}, once more
  \item \emph{Say it:} say \emph{pacci}
  \item \emph{Recall:} say the Telugu for fat, then the Telugu for thin, then say \emph{pacci} again
\end{itemize}

\begin{tcolorbox}[breakable,colback=teal!4,colframe=teal!35!black,title={Before you move on}]
What does \te{పచ్చి} mean? (\textquotedblleft{}Raw, unripe\textquotedblright{}.) Say it once more.
\end{tcolorbox}
